% !TeX program = xelatex
% ============================================================================
%  第 1 课　先导课：抽象代数是什么　—— 学生版讲义（可以写字的版本）
%  与 02-课件/第01课-先导课/第01课-先导课.tex 同步
% ============================================================================

\xhlesson{1}{先导课：抽象代数是什么}{}

\begin{mubiao}
  \begin{itemize}
    \item 知道抽象代数研究的是「运算本身的规律」，不是一批新公式。
    \item 能用集合的说法描述一堆对象（会写、会读大括号，不做集合运算）。
    \item 记住一个悬念：999999 能被 7 整除，为什么？
  \end{itemize}
\end{mubiao}

\section*{一、先记一个名字}

这门课叫《从数数开始的抽象代数生活》。名字是化用了一部动画片《Re:从零开始的异世界生活》，
不过很遗憾，这门课和那部动画没有任何关系。

如果不起一个有梗的名字，这门课只有四个字：\textbf{抽象代数}。
它是大学数学系的专业课，一般大二才开。它研究的不是「数怎么算」，而是「算」本身。

数字的世界，是怎么一点点变抽象的？看这三步：

\begin{center}
\begin{tabular}{@{}>{\raggedright\arraybackslash}p{2.4cm}>{\raggedright\arraybackslash}p{4.6cm}>{\raggedright\arraybackslash}p{8.4cm}@{}}
\toprule
\textbf{阶段} & \textbf{写出来的样子} & \textbf{这一步在干什么} \\
\midrule
小学 & \(3+5=8\) & 只用数字。 \\
初中 & \(a+b=b+a\) & 开始用字母代替数。 \\
抽象代数 & 连「运算」也用字母 & 不问「加的是谁」，只问「有什么规律」。 \\
\bottomrule
\end{tabular}
\end{center}

\noindent 这门课只有两个约定：尽量少记术语和定义；只看生活里的例子。

\section*{二、两个生活里的例子}

\begin{sikao}[例子一：体育课上的队形]
  体育课上，老师要把两列纵队变成四列纵队，他做了什么？\\[4pt]
  人还是这些人，那么变的是什么？
  \liukong{4}
\end{sikao}

\begin{sikao}[例子二：半夜 12 点]
  为什么半夜 12 点，既叫 0 点，又叫 24 点？
  \liukong{3}
\end{sikao}

\begin{example}[两个例子的共同点]
  \begin{itemize}
    \item 队形：人没变，位置上重新排了一遍。
    \item 时间：表没变，指针绕了一圈。
  \end{itemize}
  \textbf{共同点：东西没变，只是「重新排了一下」。}
\end{example}

\begin{sikao}[开放题]
  你自己再举一个例子，也是「东西没变，只是重新排了一下」：\\[4pt]
  （提示：可以想想座位表、扑克牌、跑操的队形、手机上的九宫格……）
  \liukong{4}
\end{sikao}

\section*{三、集合是什么}

\begin{definition}[集合]
  把确定的一堆东西放在一起，就叫做一个\textbf{集合}，用大括号写出来。
\end{definition}

\begin{sikao}[一个不算集合的例子]
  「我们班个子比较高的同学」，能不能打包成一个集合？\\[4pt]
  为什么？说说你的理由：
  \liukong{2}
\end{sikao}

\noindent 这个例子\textbf{不能}打包成集合——「比较高」不确定，谁算高、谁不算高，没有统一标准。
换成「我们班身高超过 \(160\) cm 的同学」，就是一个集合了：标准是确定的。
所以定义里「\textbf{确定}」两个字，是这个定义里最要紧的两个字。

\begin{example}
  \(\{1,2,3\}\)；\quad \{甲, 乙, 丙\}；\quad \{红, 黄, 蓝\}。

  所有的有理数放在一起也是一个集合；自然数、整数也都可以看作集合。
\end{example}

\noindent 这节课马上要用的三个写法：

\begin{itemize}
  \item \(\{1,2,3,\dots,12\}\) —— 钟面上的 12 个时刻；
  \item \(\{A,B,C\}\) —— 三角形的三个顶点；
  \item \{原地不动, 转 \(90^\circ\), 转 \(180^\circ\), 转 \(270^\circ\)\} —— 正方形的四种转法。
\end{itemize}

\noindent 最后一行注意：那个大括号里装的不是数，是四个\textbf{动作}。

\begin{dongshou}[自己也写两个集合]
  1. 写出一种「不是数」的东西，把它打包成一个集合：
  \liukong{2}
  2. 一个集合里可以装很多个东西，也可以只装一个。你能举出一个「只装一个东西」的集合吗？
  \liukong{2}
\end{dongshou}

在\textbf{第 12 课}，你要自己证明一个三百多年前的定理。

\section*{四、一个悬念：999999}

\[
  999999 \div 7 = \underline{\hspace{6em}}
\]

\begin{dongshou}
  先别用计算器，估一估：999999 能被 7 整除吗？在下面写下你的猜测和理由。
  \liukong{3}
\end{dongshou}

\begin{example}
  \[
    999999 \div 7 = 142857
  \]
  正好整除，一点不剩。
\end{example}

\noindent 关键线索：\(\;999999 = 1000000 - 1 = 10^{6} - 1\)。

于是问题变成：为什么 \(10^{6}-1\) 能被 7 整除？这件事和 7 有什么关系？
把 7 换成别的数还管用吗？

\begin{sikao}[把问号留在黑板上]
  这个问题今天不回答。它会一直留在黑板最右边，到第 13 课我们回来把它划掉。\\[4pt]
  在那之前，先看看下面这个「把戏」。
\end{sikao}

\begin{example}[142857 的循环]
  \begin{center}
  \begin{tabular}{@{}ll@{\hspace{2.5em}}ll@{}}
    \(142857 \times 1 = 142857\) & & \(142857 \times 4 = 571428\) & \\
    \(142857 \times 2 = 285714\) & & \(142857 \times 5 = 714285\) & \\
    \(142857 \times 3 = 428571\) & & \(142857 \times 6 = 857142\) & \\
  \end{tabular}
  \end{center}
  数字还是那六个，只是换了个位置。
\end{example}

\begin{dongshou}[自己动手算一算]
  \(142857 \times 7 = ?\)\\[6pt]
  算出来以后，和上面那行 \(999999\) 比一比，你发现了什么？
  \liukong{3}
\end{dongshou}

\section*{五、这门课要去哪}

\begin{center}
\begin{tabular}{@{}>{\raggedright\arraybackslash}p{5.4cm}>{\raggedright\arraybackslash}p{8.4cm}@{}}
\toprule
\textbf{第一站　钟表与取余} & 数着数着，会绕回来。 \\
\textbf{第二站　三角形与正方形} & 位置换一换，还是那几个样子。 \\
\textbf{第三站　换位置的算术} & 先做这个再做那个，结果会不一样。 \\
\textbf{第四站　费马小定理} & 回到今天那个 999999。 \\
\bottomrule
\end{tabular}
\end{center}

最后一节课，你会带走三样东西：

\begin{itemize}
  \item 一个能自己证明的定理（费马小定理）；
  \item 一套看世界的办法：把不同事情里相同的规矩挑出来；
  \item 一个终于能回答的问题：999999 为什么能被 7 整除。
\end{itemize}

\begin{xiaojie}
  用你自己的话说一遍，今天这节课的三件事：
  \liukong{4}
\end{xiaojie}

\noindent{\small\textbf{想读一点参考书？}下面这些书你现在大概率读不懂，至少得上高中，留个印象就好：}

\begin{itemize}
  \item 人民教育出版社．普通高中课程标准实验教科书 数学 选修 3-4 A 版：对称与群．北京：人民教育出版社，2007．
  \item 韩士安，林磊．近世代数．2 版．北京：科学出版社，2009．
  \item 内森·卡特．群论：彩图版．郭小强，罗翠玲，译．北京：机械工业出版社，2019．
  \item LEE G T. Abstract algebra: an introductory course[M]. Cham: Springer, 2018.
\end{itemize}

\noindent\textbf{下节课}：从余数到模运算——我们会给今天那个「绕圈」装上记号。